\scanpage{48}
% archive foliation: 47

Considérons cette gerbe $\unsure{\mathfrak G}$ sur $\mathbb A=\prod_p\mathbb Z_p$, de groupe $\mu_2$, définie ``au signe près'' (de façon précise, on a une $2$-catégorie des $(C,\alpha)$, qui \unsure{satisfait} $\pi_0=\pi_1=0$, $\pi_2=\mathbb Z/2\mathbb Z$, et \deleted{que} on a \deleted{\unsure{une}} \inserted{une famille de gerbes} indexée par $\mathcal D$, i.e. un $2$-foncteur de $\mathcal D$ dans la $2$-catégorie des \deleted{\unsure{$2$-}}gerbes sur $\mathbb A$, groupe $\mu_2$).
Dans cette gerbe, on a une ``\uline{section définie à isom can. au signe près}''.

\margin{attention: \unsure{réalité}, \unsure{s'il n'y a pas} \unsure{un dernier} signe!}

\medskip
\hrule
\medskip

Autre façon d'exprimer: le \deleted{catégorie des \unsure{\operatorname{Eq}}} \inserted{foncteur des} catégories de Picard \unsure{composé}
\diagram{\begin{tikzcd}[column sep=large]
\uline{\mathrm{Eq}}(\unsure{C}) \arrow[r] &
\uline{\mathrm{Eq}}(C,T) \arrow[r] &
\uline{\mathrm{Eq}}(C_p,T_p) \arrow[d,phantom,"\simeq\text{ can.}"]\\
&&\mathrm{Eq}(C(\mathbb Z_p),T(\mathbb Z_p))
\end{tikzcd}}
est pleinement fidèle, la catégorie de Picard quotient est donc une catégorie rigide, qui n'est autre que
$[0\longrightarrow\mathbb Z_p^*/\mathbb Z_p^{*2}]$.
Donc la catégorie des équivalences de $(C,T)_{\mathbb Z_p}$ avec $(C(\mathbb Z_p),T(\mathbb Z_p))$
\uline{modulo action de} $\uline{\mathrm{Eq}}(C)$ est une catégorie rigide,
\deleted{torseur} dont l'ens des classes d'isom. est un $\mathcal T_p$, torseur canonique sous $\mathbb Z_p^*/\mathbb Z_p^{*2}$ (un élément \deleted{lequel}\inserted{privilégié}!!!), indépendant à isom can. près de $C$.
\deleted{On \unsure{a un} \unsure{point} \unsure{can.} $\xi\in\mathcal T_p$.}
\deleted{Après \unsure{avoir choisi} \unsure{le dernier} \unsure{isomorphisme}}
On peut considérer les couples $(C,\alpha)$, où $\alpha$ est une équivalence
\[
(C,T)_{\mathbb Z_p}\simeq(C(\mathbb Z_p),T(\mathbb Z_p))
\]
\uline{dans la classe de} $\xi_p$, et
\deleted{\unsure{on}}
si on a deux tels couples, ils sont \uline{\unsure{isomorphes}} d'une \uuline{façon canonique}.
